\documentclass[10pt,a4paper]{amsart}
\usepackage[margin=19mm]{geometry}
\usepackage{amsmath,amssymb,mathtools}
\usepackage[scr=rsfs,cal=euler]{mathalfa}
\usepackage{xfrac}
\usepackage[unicode,hidelinks,bookmarks=false]{hyperref}
\newcommand{\ie}{i.e.\ }
\newcommand{\cf}{cf.\ }
\newcommand{\resp}{respectively\ }
\newcommand{\Subsection}{\subsection}
\providecommand{\nobreakdash}{\nobreak\hbox{-}}
\setlength{\emergencystretch}{2em}
\multlinegap0pt
\usepackage{amssymb}
\usepackage{float}
\usepackage{xcolor}
\newtheorem{thm}{Theorem}
\title{\textbf{ Quadratic form estimations for Hessian matrices of resistance distance and Kirchhoff index of positive-weighted graphs }}
\author{Yu Li, Lizhu Sun, Changjiang Bu}
\date{Source: arXiv:2603.04963v1 (CC BY 4.0)}
\begin{document}
\maketitle

\noindent\emph{Source and licence.} \textbf{ Quadratic form estimations for Hessian matrices of resistance distance and Kirchhoff index of positive-weighted graphs }. Version arXiv:2603.04963v1 (CC BY 4.0), distributed by arXiv under the Creative Commons Attribution 4.0 International licence, http://creativecommons.org/licenses/by/4.0/, as recorded in the arXiv licence metadata for this exact version and shown on the abstract page https://arxiv.org/abs/2603.04963v1. Full licence notice: \texttt{licenses/arxiv-2603.04963-CC-BY-4.0.txt}.\par\smallskip

\noindent\textit{Editorial scope.} Counted unit: the article's thm 3.3 (source0.tex lines 245--255). The article's complete original proof of it is delivered with it (source0.tex lines 256--295, one \texttt{proof} environment of 40 lines). No mathematics is written, repaired or re-derived: the statement and the proof are the article's own lines, in the article's own notation, and no retained line is substituted or re-flowed.  No further source-local context is needed: every cross-reference of every retained line resolves inside the two delivered spans.  Nothing was omitted from any retained span, so the package carries no omission record.  No retained line contains a citation, so the wrapper carries no bibliography and no bibliography entry.  No retained line contains a figure environment, an image-inclusion command or drawing code, so no image is delivered and the package declares no asset dependency.  Editorial formatting only, in addition: an AMS \texttt{amsart} preamble that restates the article's own declarations and macros used by the retained lines, namely the environments \texttt{thm} on separate counters, together with the standard packages those lines need.  The retained environments print in this document as Theorem 1 in order of appearance, whereas the article prints the same items with the numbering of its own full document.  The complete native source of the article as distributed in the arXiv e-print for this exact version is bundled unchanged as \texttt{sources/195771/source0.tex}.
\begin{thm}\label{3.3}
Let $\widetilde{A}\in\mathbb{H}^{m\times n}$, $ \widetilde{\mathbf{x}}\in\mathbb{H}^{n}$ and $\widetilde{\mathbf{b}}\in\mathbb{H}^{m}$. Suppose that $\widetilde{A}^{\dag}$ exists. Then the following statements hold.

$(a)$ The solution  to the hyper-dual  equation
$\widetilde{A}\widetilde{\mathbf{x}}=\widetilde{\mathbf{b}}$ exists  if and only if $$\widetilde{A}\widetilde{A}^{\dag}
\widetilde{\mathbf{b}}=\widetilde{\mathbf{b}}.$$

$(b)$  If the solution to the hyper-dual equation $\widetilde{A}\widetilde{\mathbf{x}}=\widetilde{\mathbf{b}}$ exists, then its general solution is given by \begin{align}\label{eq91} \widetilde{\mathbf{x}}=\widetilde{A}^{\dag}\widetilde{\mathbf{b}}
+(I-\widetilde{A}^{\dag}\widetilde{A})\widetilde{\mathbf{u}},
\end{align} where $\widetilde{\mathbf{u}}\in \mathbb{H}^{n}$ is arbitrary.
\end{thm}

\begin{proof}
$(a)$ Suppose $\widetilde{\mathbf{x}}$ is a solution to $\widetilde{A}
\widetilde{\mathbf{x}}=\widetilde{\mathbf{b}}$. Then
$$\widetilde{\mathbf{b}}=\widetilde{A}
\widetilde{\mathbf{x}}=(\widetilde{A}\widetilde{A}^{\dag})
\widetilde{A}
\widetilde{\mathbf{x}}=
\widetilde{A}\widetilde{A}^{\dag}
\widetilde{\mathbf{b}}.$$ Conversely, if
$\widetilde{A}\widetilde{A}^{\dag}
\widetilde{\mathbf{b}}
=\widetilde{\mathbf{b}}$, it is clear that $\widetilde{A}^{\dag}\widetilde{\mathbf{b}}$ is a solution to $\widetilde{A}
\widetilde{\mathbf{x}}
=\widetilde{\mathbf{b}}$.

 $(b)$  From $(a)$, we have $\widetilde{A}\widetilde{A}^{\dag}
 \widetilde{\mathbf{b}}=\widetilde{\mathbf{b}}$. For any $\widetilde{\mathbf{u}}\in \mathbb{H}^{m}$,   $$
\widetilde{A}(\widetilde{A}^{\dag}
\widetilde{\mathbf{b}}+
(I-\widetilde{A}^{\dag}\widetilde{A})
\widetilde{\mathbf{u}})
=\widetilde{A}\widetilde{A}^{\dag}\widetilde{\mathbf{b}}
=\widetilde{\mathbf{b}}.$$ Hence, $\widetilde{\mathbf{x}}=\widetilde{A}^{\dag}\widetilde{\mathbf{b}}
+(I-\widetilde{A}^{\dag}\widetilde{A})
\widetilde{\mathbf{u}}$ is a solution to $\widetilde{A}
\widetilde{\mathbf{x}}=\widetilde{\mathbf{b}}$.

On the other hand, let
$\widetilde{\mathbf{x}}$ be a solution to $\widetilde{A}
\widetilde{\mathbf{x}}=\widetilde{\mathbf{b}}$.
 Then
\begin{align}\nonumber \widetilde{\mathbf{x}}&=\widetilde{\mathbf{x}}
+\widetilde{A}^{\dag}(\widetilde{\mathbf{b}}-
\widetilde{A}\widetilde{\mathbf{x}})\\ \nonumber
&=\widetilde{A}^{\dag}\widetilde{\mathbf{b}}+
(I-\widetilde{A}^{\dag}\widetilde{A})\widetilde{\mathbf{x}}.
\end{align}
Therefore, every solution to $\widetilde{A}\widetilde{\mathbf{x}}
=\widetilde{\mathbf{b}}$ can be written in the form Eq. \eqref{eq91}.
\end{proof}

\end{document}
